\documentclass[12pt,a4paper,oneside]{article}
\usepackage[brazil]{babel}
\usepackage[utf8]{inputenc}
\usepackage[dvipsnames]{xcolor}
\usepackage{listings}
\usepackage{wrapfig}
\usepackage{olymp}

\setcounter{problem}{1}

\begin{document}

\begin{problem}{Perímetro e área do lote}{lote}{2}

\begingroup
\setlength{\intextsep}{0pt}%
%\setlength{\columnsep}{0pt}

José comprou um lote retangular e agora precisa estimar os gastos para cercar e limpar seu lote.
Para murá-lo, ele precisa saber o perímetro. Para capiná-lo, precisa saber a área.
Ajude José com um programa que calcula o perímetro e a área, dados o comprimento e a largura do lote.

Por exemplo, se o lote tem $20$ de largura e $30$ de comprimento, seu perímetro é $20+20+30+30 = 100$ e sua área é $20 \times 30 = 600$

% \Notes

\InputFile

A entrada contém apenas uma linha, com dois valores inteiros, o comprimento $C$ e a largura $L$ do lote.

Restrições: $1 \leq C, L \leq 1000$.

\endgroup

\OutputFile

Escreva uma linha contendo o perímetro e a área do lote.

% \Notes

\Example

\begin{example}
\exmp{
20 30
}{
100 600
}%
\end{example}

\begin{example}
\exmp{
12 20
}{
64 240
}%
\end{example}

\begin{example}
\exmp{
100 100
}{
400 10000
}%
\end{example}

%Comentários sobre o exemplo, se necessário.

\end{problem}

\end{document}